\section{Problem 2: Simulation and Analysis of Stock Price Trajectories}
\subsection{Description}
Simulate five sample trajectories of \eqref{c1:model} for the following parameter values and plot the trajectories on the same set of coordinate axes: $\mu = 0.12, \sigma = 0.1, T = 1, s = \$40, N = 254$. Then repeat the experiment using the value $\sigma = 0.25$ for the volatility. Do the sample trajectories generated in the latter case appear to exhibit a greater degree of variability in their behavior?

\textit{Hint:} For the $\varepsilon_n$'s, it is permissible to use a random number generator that creates normally distributed random numbers with mean zero and variance $1$.

\subsection{Simulation}
To analyze the effect of volatility on stock price behavior, we simulated five sample trajectories using the discrete model
\begin{equation*}
	S_{n+1} = S_n + \mu S_n \Delta t + \sigma S_n \varepsilon_{n+1} \sqrt{\Delta t}.
\end{equation*}
Five trajectories with different colors were generated with parameters $\mu = 0.12$, $T = 1$, $s = \$40$ and $N = 254$.

\begin{figure}[ht!]
	\centering
	\begin{tikzpicture}
		\begin{axis}[
			width=0.48\textwidth,
			title={Volatility $\sigma = 0.1$},
			xlabel={Time $t$},
			ylabel={Stock Price $S_n$},
			grid=both,
			minor tick num=1,
			legend style={font=\tiny},
			ymax=55, ymin=35
			]
			\addplot[blue, thick] table [x=Time, y=Low_0, col sep=comma] {trajectories.csv};
			\addplot[red, thick]  table [x=Time, y=Low_1, col sep=comma] {trajectories.csv};
			\addplot[green, thick] table [x=Time, y=Low_2, col sep=comma] {trajectories.csv};
			\addplot[orange, thick] table [x=Time, y=Low_3, col sep=comma] {trajectories.csv};
			\addplot[cyan, thick] table [x=Time, y=Low_4, col sep=comma] {trajectories.csv};
		\end{axis}
	\end{tikzpicture}
	\hfill
	\begin{tikzpicture}
		\begin{axis}[
			width=0.48\textwidth,
			title={Volatility $\sigma = 0.25$},
			xlabel={Time $t$},
			grid=both,
			minor tick num=1,
			ymax=55, ymin=35
			]
			\addplot[blue, thick] table [x=Time, y=High_0, col sep=comma] {trajectories.csv};
			\addplot[red, thick]  table [x=Time, y=High_1, col sep=comma] {trajectories.csv};
			\addplot[green, thick] table [x=Time, y=High_2, col sep=comma] {trajectories.csv};
			\addplot[orange, thick] table [x=Time, y=High_3, col sep=comma] {trajectories.csv};
			\addplot[cyan, thick] table [x=Time, y=High_4, col sep=comma] {trajectories.csv};
		\end{axis}
	\end{tikzpicture}
	\caption{Comparison of stock price trajectories for low and high volatility.}
	\label{fig:trajectories}
\end{figure}

\subsection{Analysis}
As observed in \figref{fig:trajectories}, the sample trajectories generated with $\sigma = 0.25$ exhibit a significantly greater degree of variability.

While the paths for $\sigma = 0.1$ remain relatively close to the deterministic growth trend, the paths for $\sigma = 0.25$ show large, erratic fluctuations and a much wider spread at the expiration time $T=1$.

This confirms that $\sigma$ serves as a scaling factor for the stochastic shocks in the model.